% Convenzioni: c=hbar=k_B=1 se non diversamente indicato; G esplicita.

\section{Nane bianche: materia degenere ed equilibrio}
\subsection{Quadro fisico e storico}

% 1
\begin{frame}[t]{Capitolo 5: struttura e obiettivi}
\small
La nana bianca è il residuo stellare il cui sostegno alla gravità è dato soprattutto dalla \emph{pressione di degenerazione degli elettroni}. I nuclei determinano quasi tutta la massa.
\begin{enumerate}
\item Dalle politropiche alla relazione massa--raggio (§5.2).
\item Equazione differenziale e soluzioni di Lane--Emden (§5.3).
\item Limite di Chandrasekhar (§5.4).
\item Correzioni coulombiane e strati della nana bianca (§§5.5--5.6).
\item Degenerazione termica, cristallizzazione e osservazioni (§§5.7--5.8).
\end{enumerate}
Convenzioni: unità naturali $c=\hbar=k_B=1$ se non esplicitato; $\rho$ è densità di \emph{massa} e $\varepsilon$ densità di energia. Nei calcoli newtoniani $\rho\simeq\varepsilon/c^2$.
\end{frame}

% 2
\begin{frame}[t]{§5.1: la scoperta e la spiegazione delle nane bianche}
\small
Sirius B fu osservata nel 1862 come compagna di Sirius A. Nonostante una massa dell'ordine della massa solare, ha un raggio confrontabile con quello terrestre.
\[
\bar\rho=\frac{3M}{4\pi R^3}\gg\rho_{\odot}.
\]
\begin{itemize}
\item Il redshift gravitazionale delle righe spettrali conferma l'elevata compattezza.
\item Fowler (1926) applica la statistica di Fermi--Dirac alla pressione degli elettroni.
\item Chandrasekhar include la relatività speciale e trova una massa limite.
\end{itemize}
Valori storici citati dal libro per Sirius B: $M\sim1.02M_\odot$, $R\sim0.0086R_\odot$, $T_{\rm eff}\sim2.5\times10^4\,{\rm K}$. I dati successivi della Tabella 5.3 sono leggermente differenti.
\end{frame}

% 3
\begin{frame}[t]{Il sostegno gravitazionale: chi fornisce massa e pressione?}
\small
Materia composta da nuclei $(A,Z)$ ed elettroni quasi liberi, in neutralità locale:
\[
n_e=Z n_A,\qquad n_B=A n_A=\frac AZ n_e,\qquad
\rho\simeq A m_u n_A=\frac AZ m_u n_e.
\]
$m_u$ è l'unità di massa atomica; $\mu_e\equiv A/Z$ è il peso molecolare medio per elettrone (\emph{non} il potenziale chimico degli elettroni, indicato $\mu_{e,\rm chem}$ quando necessario).
\[
P\simeq P_e(n_e),\qquad \varepsilon\simeq\rho c^2
\]
finché nuclei non relativistici ed elettroni degeneri dominano, con effetti coulombiani, termici e composizionali come correzioni.
\end{frame}

\subsection{§5.2: politropiche e relazione massa--raggio}

% 4
\begin{frame}[t]{Polytropic EOS: dalla densità numerica alla densità di massa}
\small
Il libro usa, a bassa energia, una legge di potenza
\[
P=K n^\Gamma,\qquad\rho=m n,\qquad
P=K'\rho^\Gamma,\quad K'=K/m^\Gamma.
\]
$\Gamma$ è l'esponente politropico della densità, $K,K'$ costanti per un dato modello. Non si deve identificare automaticamente $\Gamma$ con l'esponente della relazione $P(\varepsilon)$ a densità ultrarelativistiche.
\medskip
Per una nana bianca, $n$ può essere $n_e$, ma allora la conversione corretta è $\rho\simeq\mu_e m_u n_e$, non $m_e n_e$.
\medskip
Relazioni già ottenute dal capitolo 3:
\[
P_e^{\rm NR}\propto n_e^{5/3},\qquad P_e^{\rm UR}\propto n_e^{4/3}.
\]
\end{frame}

% 5
\begin{frame}[t]{Equilibrio idrostatico newtoniano: due equazioni}
\small
Per simmetria sferica la massa interna a $r$ è
\[
m(r)=4\pi\int_0^r r'^2\rho(r')\,dr',\qquad
\frac{dm}{dr}=4\pi r^2\rho.
\]
Il bilancio fra forza di gravità e gradiente di pressione richiede
\[
\frac{dP}{dr}=-\rho(r)\frac{Gm(r)}{r^2}.
\]
Dati $P(\rho)$ e $\rho_c$, integriamo da $m(0)=0$ e $P(0)=P_c$ fino al primo zero di $P$:
\[
P(R)=0,\qquad M=m(R).
\]
È l'approssimazione newtoniana appropriata alla gran parte delle nane bianche; la relatività speciale entra già nell'EOS degli elettroni.
\end{frame}

% 6
\begin{frame}[t]{Stima del gradiente e della pressione centrale}
\small
Una stima dimensionale sostituisce il profilo dettagliato con
\[
\left|\frac{dP}{dr}\right|\sim\frac{P_c}{R},\quad
m(R)=M,\quad \bar\rho=\frac{3M}{4\pi R^3}.
\]
Dall'equilibrio:
\[
\frac{P_c}{R}\sim\frac{GM\bar\rho}{R^2},\qquad
P_c\sim\frac{GM\bar\rho}{R}\sim\frac{GM^2}{R^4}.
\]
Un coefficiente geometrico dell'ordine dell'unità è stato tralasciato: questa non è una soluzione esatta e $\rho_c\neq\bar\rho$ in generale.
\medskip
La relazione mette in evidenza la competizione fra gravità ($GM^2/R^4$) e pressione microfisica.
\end{frame}

% 7
\begin{frame}[t]{Eliminare la pressione: la legge massa--raggio}
\small
Per una politropica $P=K'\rho^\Gamma$, sostituiamo $\rho\sim M/R^3$ nella stima:
\[
K'\left(\frac{M}{R^3}\right)^\Gamma\sim \frac{GM^2}{R^4}.
\]
Separando le potenze:
\[
K' M^\Gamma R^{-3\Gamma}\sim GM^2R^{-4},
\]
\[
\frac{K'}G\sim M^{2-\Gamma}R^{3\Gamma-4}.
\]
Dunque, per $K'$ fissato,
\[
M^{2-\Gamma}R^{3\Gamma-4}\simeq\mathrm{costante}.
\]
Questa è l'equazione (5.10), a meno del coefficiente senza dimensioni che sarà determinato mediante Lane--Emden.
\end{frame}

% 8
\begin{frame}[t]{Indice politropico e stabilità: $\Gamma=4/3$}
\small
Nel modello newtoniano omologo, compressione a massa fissata implica $\rho\sim R^{-3}$:
\[
P_{\rm materia}\propto R^{-3\Gamma},\qquad
P_{\rm grav}\propto R^{-4}.
\]
Se $\Gamma<4/3$, la pressione gravitazionale cresce più rapidamente, favorendo l'instabilità a una compressione. Per $\Gamma>4/3$ il termine materiale può ristabilire l'equilibrio.
\medskip
Al limite $\Gamma=4/3$:
\[
M^{2/3}\propto K'/G\quad\Longrightarrow\quad
M\propto(K'/G)^{3/2},
\]
indipendentemente da $R$: è l'origine dimensionale di $M_{\rm Ch}$.
\medskip
\emph{Precisazione}: $\Gamma=4/3$ è una soglia del problema newtoniano idealizzato; non garantisce stabilità completa né descrive l'esito finale di un collasso.
\end{frame}

% 9
\begin{frame}[t]{Limite non relativistico: $\Gamma=5/3$}
\small
Per elettroni non relativistici degeneri:
\[
P_e^{\rm NR}=K'_{\rm NR}\rho^{5/3}.
\]
Inseriamo $\Gamma=5/3$ nella relazione
\[
M^{2-\Gamma}R^{3\Gamma-4}=\mathrm{cost.}:
\]
\[
M^{1/3}R^1=\mathrm{cost.}\quad\Longrightarrow\quad
MR^3=\mathrm{cost.}
\]
Quindi, a composizione fissa, aumentando la massa il raggio diminuisce:
\[
R\propto M^{-1/3}.
\]
Questo descrive il ramo di nane bianche non relativistiche, lontano dal limite di Chandrasekhar.
\end{frame}

% 10
\begin{frame}[t]{Altri indici: $\Gamma=2$ e fluido incomprimibile}
\small
Per $\Gamma=2$:
\[
M^0R^2\propto K'/G\quad\Longrightarrow\quad R=\mathrm{cost.}
\]
Un esempio modellistico è un'EOS dominata da interazioni repulsive proporzionali a $n^2$.
\medskip
Nel limite $\Gamma\to\infty$, per una materia a densità fissata,
\[
M=\frac{4\pi}3\rho R^3\quad\Longrightarrow\quad M\propto R^3.
\]
Una sfera incomprimibile ha densità spazialmente costante, ma ciò \emph{non implica da solo} che sia autolegata: occorre una microfisica che mantenga densità non nulla anche a $P=0$.
\end{frame}

\section{L'equazione di Lane--Emden}
\subsection{§5.3: derivazione e soluzioni}

% 11
\begin{frame}[t]{Eliminare la massa $m(r)$ dall'equilibrio}
\small
Partiamo da
\[
\frac{dP}{dr}=-\rho\frac{Gm}{r^2},\qquad
\frac{dm}{dr}=4\pi r^2\rho.
\]
Moltiplichiamo la prima per $r^2/\rho$:
\[
\frac{r^2}{\rho}\frac{dP}{dr}=-Gm(r).
\]
Deriviamo rispetto a $r$:
\[
\frac{d}{dr}\left(\frac{r^2}{\rho}\frac{dP}{dr}\right)
=-G\frac{dm}{dr}=-4\pi G r^2\rho.
\]
Dopo divisione per $r^2$:
\[
\frac1{r^2}\frac{d}{dr}\left(\frac{r^2}{\rho}\frac{dP}{dr}\right)=-4\pi G\rho.
\]
È l'equazione (5.12) del testo.
\end{frame}

% 12
\begin{frame}[t]{Collegamento corretto con il potenziale newtoniano}
\small
Definiamo l'accelerazione gravitazionale tramite $-\nabla\Phi$. In simmetria sferica:
\[
\frac{d\Phi}{dr}=\frac{Gm(r)}{r^2},\qquad
\frac{dP}{dr}=-\rho\frac{d\Phi}{dr}.
\]
L'equazione di Poisson è
\[
\nabla^2\Phi=\frac1{r^2}\frac{d}{dr}
\left(r^2\frac{d\Phi}{dr}\right)=4\pi G\rho.
\]
Eliminando $\Phi$ riotteniamo
\[
\frac1{r^2}\frac{d}{dr}\left(\frac{r^2}{\rho}\frac{dP}{dr}\right)=-4\pi G\rho.
\]
\emph{Nota di revisione}: con $\Phi=-GM/r$ all'esterno, il segno di Poisson è \emph{positivo}; le equazioni (5.13)--(5.14) nel testo, così come stampate, hanno una discordanza di segno rispetto alla convenzione abituale.
\end{frame}

% 13
\begin{frame}[t]{Dall'esponente $\Gamma$ all'indice $n$ di Lane--Emden}
\small
Definiamo il polytropic index (non la densità numerica!) come
\[
n=\frac1{\Gamma-1},\qquad\Gamma=1+\frac1n.
\]
Introduciamo le variabili senza dimensioni:
\[
r=a\xi,\qquad \rho(r)=\rho_c\theta(\xi)^n,\qquad \theta(0)=1.
\]
Con $P=K'\rho^{1+1/n}$ segue
\[
P(r)=K'\rho_c^{1+1/n}\theta^{n+1}.
\]
Il parametro $a$ è una scala di lunghezza; $\xi$ misura la distanza in unità di $a$, mentre $\theta$ rappresenta la radice normalizzata del profilo di densità.
\end{frame}

% 14
\begin{frame}[t]{Derivata della pressione nella nuova variabile}
\small
Da $\rho=\rho_c\theta^n$:
\[
\frac{d\rho}{d\xi}=n\rho_c\theta^{n-1}\theta',
\quad \frac{dP}{d\rho}=K'\left(1+\frac1n\right)\rho^{1/n}.
\]
Il prodotto semplifica:
\[
\frac1\rho\frac{dP}{dr}
=\frac{(n+1)K'\rho_c^{1/n}}{a}\,\theta'(\xi).
\]
Dunque
\[
\frac{r^2}{\rho}\frac{dP}{dr}
=(n+1)K'\rho_c^{1/n}a\,\xi^2\theta'(\xi).
\]
L'idea della sostituzione è trasformare un'espressione non lineare in $\rho$ in una derivata semplice di $\theta$.
\end{frame}

% 15
\begin{frame}[t]{Scelta della scala $a$ e forma di Lane--Emden}
\small
Inseriamo la sostituzione in
\[
\frac1{r^2}\frac{d}{dr}\left(\frac{r^2}{\rho}\frac{dP}{dr}\right)=-4\pi G\rho_c\theta^n.
\]
Segue
\[
\frac{(n+1)K'\rho_c^{1/n-1}}{4\pi G a^2}
\frac1{\xi^2}\frac d{d\xi}(\xi^2\theta')=-\theta^n.
\]
Scegliendo
\[
a^2=\frac{(n+1)K'\rho_c^{1/n-1}}{4\pi G}
\]
il coefficiente diventa 1 e otteniamo
\[
\frac1{\xi^2}\frac d{d\xi}
\left(\xi^2\frac{d\theta}{d\xi}\right)=-\theta^n.
\]
\end{frame}

% 16
\begin{frame}[t]{Condizioni al centro e all'esterno}
\small
La soluzione fisica regolare al centro soddisfa
\[
\theta(0)=1,\qquad \theta'(0)=0.
\]
La prima è una normalizzazione a $\rho_c$; la seconda è imposta da regolarità e simmetria sferica.
\medskip
La superficie, per una politropica con densità tendente a zero, corrisponde al primo zero positivo:
\[
\theta(\xi_1)=0,\qquad R=a\xi_1.
\]
Il problema differenziale è singolare in $\xi=0$ se scritto $\theta''+2\theta'/\xi=-\theta^n$, ma la soluzione fisica è regolare.
\medskip
Per $n<5$ essa ha un primo zero a raggio finito; per $n=5$ il raggio è infinito, ma la massa totale resta finita.
\end{frame}

% 17
\begin{frame}[t]{Sviluppo in serie della soluzione presso il centro}
\small
Poniamo
\[
\theta(\xi)=1+a_2\xi^2+a_4\xi^4+\cdots,\qquad
\theta'=2a_2\xi+4a_4\xi^3+\cdots.
\]
L'equazione di Lane--Emden dà
\[
\theta''+\frac{2}{\xi}\theta'
=6a_2+20a_4\xi^2+\cdots=-\theta^n.
\]
Espandendo $\theta^n=1+n a_2\xi^2+\cdots$:
\[
6a_2=-1,\qquad20a_4=-na_2.
\]
Perciò
\[
\theta(\xi)=1-\frac{\xi^2}{6}+\frac{n\xi^4}{120}+O(\xi^6).
\]
Questa serie permette di avviare un'integrazione numerica stabile a $\xi>0$ piccolo.
\end{frame}

% 18
\begin{frame}[t]{Tre soluzioni analitiche di Lane--Emden}
\small
L'equazione ammette soluzioni elementari per tre indici:
\[
\begin{array}{c|c|c}
n & \theta(\xi) & \text{primo zero }\xi_1\\ \hline
0 & 1-\xi^2/6 & \sqrt6\\[1mm]
1 & \sin\xi/\xi & \pi\\[1mm]
5 & (1+\xi^2/3)^{-1/2} & \infty
\end{array}
\]
Per $n=0$ la densità è costante (sfera incomprimibile); per $n=1$, $\Gamma=2$ e il raggio è indipendente da $\rho_c$.
\medskip
Nel caso $n=5$, $\rho\propto\theta^5$ decresce sufficientemente in fretta da produrre massa totale finita, anche se il raggio formale diverge.
\end{frame}

% 19
\begin{frame}[t]{Profilo di Lane--Emden: confronto fra diversi $n$}
\centering
\begin{figure}[!ht]\includegraphics[trim={0cm 0cm 0 0},clip,keepaspectratio,width=0.95\textwidth]{chatGPT/ch5_lane_emden_curve}\label{fig:chatGPT_ch5_lane_emden_curve}\end{figure}
\par\scriptsize Grafico calcolato numericamente dagli appunti con $\theta(0)=1,\theta'(0)=0$.
\end{frame}

% 20
\begin{frame}[t]{Massa senza integrali: l'identità di Lane--Emden}
\small
Dal profilo $\rho=\rho_c\theta^n$:
\[
M=4\pi\int_0^R r^2\rho(r)\,dr
=4\pi a^3\rho_c\int_0^{\xi_1}\xi^2\theta^n\,d\xi.
\]
Ma dall'equazione di Lane--Emden,
\[
\xi^2\theta^n=-\frac{d}{d\xi}(\xi^2\theta').
\]
Integrando:
\[
\int_0^{\xi_1}\xi^2\theta^n d\xi
=-\left[\xi^2\theta'\right]_0^{\xi_1}
=-\xi_1^2\theta'(\xi_1).
\]
Definiamo $\omega_n\equiv-\xi_1^2\theta'(\xi_1)>0$:
\[
M=4\pi a^3\rho_c\omega_n.
\]
\end{frame}

% 21
\begin{frame}[t]{Leggi esatte di scala di $R$ e $M$}
\small
Ricordiamo
\[
a^2=\frac{(n+1)K'}{4\pi G}\rho_c^{(1-n)/n}.
\]
Poiché $R=\xi_1a$, otteniamo
\[
R=\xi_1\left[\frac{(n+1)K'}{4\pi G}\right]^{1/2}
\rho_c^{(1-n)/(2n)}.
\]
Da $M=4\pi\omega_n a^3\rho_c$ segue
\[
M=4\pi\omega_n
\left[\frac{(n+1)K'}{4\pi G}\right]^{3/2}
\rho_c^{(3-n)/(2n)}.
\]
A EOS fissata, le soluzioni differiscono solo per la densità centrale $\rho_c$.
\end{frame}

% 22
\begin{frame}[t]{Da Lane--Emden alla legge massa--raggio}
\small
La scala del raggio dà $R\propto\rho_c^{(1-n)/(2n)}$; la massa dà $M\propto\rho_c^{(3-n)/(2n)}$.
\medskip
Eliminando $\rho_c$:
\[
M^{n-1}R^{3-n}=\mathrm{costante}(K',G,n).
\]
Con $\Gamma=1+1/n$ otteniamo
\[
M^{2-\Gamma}R^{3\Gamma-4}=\mathrm{costante}.
\]
La stessa legge del ragionamento dimensionale §5.2 emerge ora \emph{esattamente}, con coefficiente numerico determinato da $\xi_1$ e $\omega_n$.
\medskip
L'argomento è strettamente newtoniano e richiede una singola politropica valida nell'intera configurazione.
\end{frame}

\section{Massa di Chandrasekhar}
\subsection{§5.4: elettroni degeneri ultrarelativistici}

% 23
\begin{frame}[t]{Quali indici servono per le nane bianche?}
\small
\[
\begin{array}{c|c|c}
\text{elettroni} & \Gamma & n=1/(\Gamma-1)\\ \hline
\text{non relativistici} &5/3&3/2\\
\text{ultrarelativistici}&4/3&3
\end{array}
\]
Per questi valori l'equazione di Lane--Emden non ammette soluzione elementare.
\medskip
Valori riportati nella Tabella 5.1 del libro:
\[
\begin{array}{c|c|c}
n& \xi_1 & \omega_n=-\xi_1^2\theta'(\xi_1)\\\hline
3/2&3.65375&2.71406\\
3&6.89685&2.01824
\end{array}
\]
La nostra integrazione numerica riproduce i valori con le cifre indicate.
\end{frame}

% 24
\begin{frame}[t]{Pressione elettronica nel limite ultrarelativistico}
\small
Per spin $1/2$ e $T=0$, $g=2$:
\[
n_e=2\int_{|\mathbf k|<k_F}\frac{d^3k}{(2\pi\hbar)^3}
=\frac{k_F^3}{3\pi^2\hbar^3}.
\]
Nel limite $E\simeq pc$ e $v\simeq c$:
\[
P_e=\frac{2}{3}\int^{k_F}\frac{d^3p}{(2\pi\hbar)^3}pv
=\frac{c\,k_F^4}{12\pi^2\hbar^3}.
\]
Invertendo $k_F=\hbar(3\pi^2n_e)^{1/3}$:
\[
P_e=\frac{(3\pi^2)^{1/3}}4\hbar c\,n_e^{4/3}.
\]
È la politropica con $\Gamma=4/3$.
\end{frame}

% 25
\begin{frame}[t]{Da $n_e$ a $\rho$: la costante $K'_{\rm UR}$}
\small
Neutralità elettrica e massa dominata dai nuclei:
\[
n_e=\frac{Z}{A}\frac\rho{m_u}=\frac\rho{\mu_e m_u},\qquad\mu_e\equiv\frac AZ.
\]
Quindi
\[
P=K'_{\rm UR}\rho^{4/3},\qquad
K'_{\rm UR}=\frac{(3\pi^2)^{1/3}}4\frac{\hbar c}{(\mu_e m_u)^{4/3}}.
\]
Non è la massa dell'elettrone a determinare la massa limite: $m_e$ si elimina nel regime ultrarelativistico, mentre resta la massa per elettrone $\mu_e m_u$.
\end{frame}

% 26
\begin{frame}[t]{Il limite $n=3$: $M$ non dipende da $\rho_c$}
\small
Dalla formula generale
\[
M=4\pi\omega_n\left[\frac{(n+1)K'}{4\pi G}\right]^{3/2}
\rho_c^{(3-n)/(2n)},
\]
ponendo $n=3$ troviamo l'esponente nullo:
\[
M=4\pi\omega_3\left(\frac{K'_{\rm UR}}{\pi G}\right)^{3/2}.
\]
Il raggio, invece, scala come
\[
R\propto\rho_c^{-1/3}.
\]
Il modello limite ultrarelativistico ideale produce quindi una massa indipendente dalla densità centrale: è il limite di Chandrasekhar, non una sequenza fisica di stelle con ogni raggio possibile.
\end{frame}

% 27
\begin{frame}[t]{Derivazione completa del coefficiente di Chandrasekhar}
\small
Inseriamo $K'_{\rm UR}=(3\pi^2)^{1/3}\hbar c/[4(\mu_e m_u)^{4/3}]$:
\[
M_{\rm Ch}=4\pi\omega_3\left[
\frac{(3\pi^2)^{1/3}\hbar c}{4\pi G(\mu_e m_u)^{4/3}}
\right]^{3/2}.
\]
Semplificando le potenze:
\[
M_{\rm Ch}=\frac{\sqrt{3\pi}}2\,\omega_3
\frac{(\hbar c/G)^{3/2}}{\mu_e^2m_u^2}.
\]
Con $\omega_3=2.01824$ e $m_{\rm Pl}=\sqrt{\hbar c/G}$:
\[
M_{\rm Ch}=3.098\frac{m_{\rm Pl}^3}{\mu_e^2m_u^2}
\simeq\frac{5.83}{\mu_e^2}M_\odot.
\]
Per $\mu_e=2$ (He, C, O): $M_{\rm Ch}\simeq1.46M_\odot$.
\end{frame}

% 28
\begin{frame}[t]{Dipendenza dalla composizione nucleare}
\small
Il libro riassume la dipendenza nell'equazione (5.40):
\[
M_{\rm Ch}\simeq1.456M_\odot
\left(\frac{2Z}{A}\right)^2\left(\frac{m_u}{m_A/A}\right)^2.
\]
Per nuclei $^4\mathrm{He}$, $^{12}\mathrm C$, $^{16}\mathrm O$, $Z/A=1/2$; per $^{56}\mathrm{Fe}$, $Z/A=26/56\simeq0.464$.
\medskip
Mantenendo uguale il resto:
\[
\frac{M_{\rm Ch}(\mathrm{Fe})}{M_{\rm Ch}(Z/A=1/2)}
\sim\left(\frac{26/56}{1/2}\right)^2\simeq0.86.
\]
Il libro riporta $M_{\rm Ch}(\mathrm{Fe})\simeq1.26M_\odot$. Il nucleo determina l'inerzia gravitazionale ma non la pressione elettronica, perciò la frazione $Z/A$ è cruciale.
\end{frame}

% 29
\begin{frame}[t]{Perché $M_{\rm Ch}$ non è la massa di tutte le nane bianche?}
\small
La deduzione usa assunzioni ideali:
\begin{itemize}
\item gas elettronico libero e completamente degenere;
\item limite \emph{ultrarelativistico} per gli elettroni;
\item nuclei non relativistici, composizione fissa e temperatura trascurata;
\item equilibrio newtoniano, senza rotazione né campi magnetici forti;
\item nessuna cattura elettronica o reazione nucleare destabilizzante.
\end{itemize}
Una nana bianca ordinaria con elettroni solo parzialmente relativistici segue una curva $M(R)$ con $R$ decrescente all'aumentare della massa; gli effetti addizionali possono spostare la massa critica rispetto al valore ideale.
\end{frame}

\section{Correzioni coulombiane}
\subsection{§5.5: cella di Wigner--Seitz}

% 30
\begin{frame}[t]{Perché la repulsione coulombiana modifica l'EOS?}
\small
I nuclei carichi positivamente formano un reticolo immerso in elettroni degeneri. Per stimare il contributo elettrostatico, approssimiamo ogni regione con una cella sferica di Wigner--Seitz:
\[
V_{\rm cell}=\frac{4\pi}3a_i^3,\qquad n_i=V_{\rm cell}^{-1}.
\]
Neutralità della singola cella:
\[
Z=n_eV_{\rm cell}=n_e\frac{4\pi}3a_i^3,
\qquad a_i=\left(\frac{3Z}{4\pi n_e}\right)^{1/3}.
\]
Indicheremo $e^2\equiv\alpha\hbar c$ in unità gaussiane per il termine $e^2/r$, con $\alpha\simeq1/137$. Questa convenzione evita di mescolare con il sistema razionalizzato $e^2/(4\pi r)$.
\end{frame}

% 31
\begin{frame}[t]{Energia elettrostatica della nube elettronica}
\small
La carica totale degli elettroni è $-Ze$; supponendoli uniformemente distribuiti nella sfera di raggio $a_i$, la carica interna a $r$ è
\[
Q_e(r)=-Ze\left(\frac r{a_i}\right)^3.
\]
L'energia di auto-interazione elettrostatica di una sfera uniformemente carica è
\[
E_{ee}=\frac35\frac{Z^2e^2}{a_i}.
\]
Si ottiene integrando l'energia dei gusci assemblati dal nulla, $dE=Q(r)\,dQ(r)/r$ (a convenzione elettrostatica fissata).
\medskip
Il termine è positivo: cariche dello stesso segno si respingono.
\end{frame}

% 32
\begin{frame}[t]{Energia attrattiva ione--elettroni}
\small
Il nucleo, trattato come carica puntiforme $+Ze$, interagisce con la nube uniforme:
\[
\rho_e=-\frac{3Ze}{4\pi a_i^3},\qquad dQ_e=4\pi r^2\rho_e\,dr.
\]
L'energia di interazione è
\[
E_{ie}=\int_0^{a_i}\frac{Ze\,dQ_e}{r}
=4\pi Ze\rho_e\int_0^{a_i}r\,dr.
\]
Sostituendo $\rho_e$ e svolgendo l'integrale:
\[
E_{ie}=-\frac32\frac{Z^2e^2}{a_i}.
\]
Il segno è negativo, perché attrazione fra cariche opposte.
\end{frame}

% 33
\begin{frame}[t]{Somma coulombiana e confronto con il reticolo reale}
\small
\[
E_C=E_{ee}+E_{ie}
=\left(\frac35-\frac32\right)\frac{Z^2e^2}{a_i}
=-\frac9{10}\frac{Z^2e^2}{a_i}.
\]
Il coefficiente $0.900$ ottenuto con la cella sferica è vicino a quello di una rete cubica a corpo centrato (bcc): $0.90981$ per la corrispondente energia di Madelung.
\medskip
La correzione è negativa: la disposizione ordinata e neutra riduce l'energia rispetto ai nuclei e agli elettroni separati.
\medskip
In questa stima non compaiono ancora effetti di scambio, correlazioni elettroniche e deformazioni della nube.
\end{frame}

% 34
\begin{frame}[t]{Energia coulombiana per elettrone}
\small
Usiamo la neutralità $a_i^{-1}=(4\pi n_e/3Z)^{1/3}$:
\[
\frac{E_C}Z=-\frac9{10}Z^{2/3}e^2
\left(\frac{4\pi n_e}{3}\right)^{1/3}.
\]
Ponendo $A_C\equiv\frac9{10}Z^{2/3}e^2(4\pi/3)^{1/3}>0$,
\[
e_C(n_e)\equiv\frac{E_C}Z=-A_C n_e^{1/3}.
\]
L'energia di volume del reticolo è
\[
\varepsilon_C=n_ee_C=-A_C n_e^{4/3}<0.
\]
Quindi la correzione presenta la stessa potenza della densità degli elettroni ultrarelativistici, ma segno opposto.
\end{frame}

% 35
\begin{frame}[t]{Pressione coulombiana e identità $P_C=\varepsilon_C/3$}
\small
A temperatura zero, la pressione associata all'energia per elettrone è
\[
P_C=n_e^2\frac{de_C}{dn_e}.
\]
Da $e_C=-A_Cn_e^{1/3}$:
\[
\frac{de_C}{dn_e}=-\frac{A_C}{3}n_e^{-2/3},\qquad
P_C=-\frac{A_C}{3}n_e^{4/3}.
\]
Poiché $\varepsilon_C=-A_Cn_e^{4/3}$,
\[
P_C=\frac{\varepsilon_C}{3}<0.
\]
La struttura del reticolo \emph{riduce} la pressione totale, senza cambiare l'esponente $4/3$ della correzione.
\end{frame}

% 36
\begin{frame}[t]{Quanto conta la correzione coulombiana?}
\small
Per elettroni ultrarelativistici,
\[
P_e\propto \hbar c\,n_e^{4/3},\qquad
|P_C|\propto Z^{2/3}\alpha\hbar c\,n_e^{4/3}.
\]
La frazione correttiva è quindi dell'ordine
\[
\frac{|P_C|}{P_e}\sim \alpha Z^{2/3}.
\]
Il libro indica una scala dell'ordine dell'1\% per He e del 3\% per O.
\medskip
A densità molto basse $P_e^{\rm NR}\propto n_e^{5/3}$ decresce più rapidamente di $|P_C|\propto n_e^{4/3}$; i due termini possono diventare confrontabili. In quel regime il modello di gas elettronico libero e completamente ionizzato va però trattato con cautela.
\end{frame}

% 37
\begin{frame}[t]{Transizione da stella gravitazionale a materia autolegata?}
\small
Il libro usa il confronto fra $P_e>0$ e $P_C<0$ per stimare una densità critica molto bassa (nell'intervallo di pochi--decine di $\mathrm{g\,cm^{-3}}$) per la quale
\[
P_{\rm totale}\simeq P_e+P_C=0.
\]
L'idea fisica è che una struttura solida può possedere un equilibrio a densità finita senza la compressione gravitazionale.
\medskip
Nel limite di densità quasi costante,
\[
M\simeq\frac{4\pi}3\rho R^3
\quad\Longrightarrow\quad M\propto R^3.
\]
\emph{Avvertenza}: una superficie a $P=0$ deriva qui da una microfisica di legame elettrico, non dalla sola approssimazione incomprimibile. Alle densità planetarie occorrono anche struttura atomica, ionizzazione parziale e proprietà dei solidi.
\end{frame}

\section{Struttura interna delle nane bianche}
\subsection{§5.6: soglie di densità e spessore degli strati}

% 38
\begin{frame}[t]{Dalla superficie verso il centro}
\small
Lo schema del testo (da usare come modello di massima):
\begin{enumerate}
\item \textbf{Atmosfera}: elettroni legati agli atomi, spettri di H o He.
\item \textbf{Interno ionizzato}: elettroni degeneri e nuclei; secondo $T$ e densità, nuclei in reticolo o fluido Coulombiano (\emph{ocean}).
\item \textbf{Zone interne dense}: elettroni possono diventare relativistici, modificando $\Gamma$ da $5/3$ verso $4/3$.
\item \textbf{Core vicino al limite}: possibili catture elettroniche e cambiamenti di composizione, dipendenti dalle soglie nucleari.
\end{enumerate}
A differenza di una stella di neutroni, il sostegno principale proviene dagli elettroni, mentre la massa è nucleonica.
\end{frame}

% 39
\begin{frame}[t]{Relazione fra densità, composizione e momento di Fermi}
\small
Per elettroni spin $1/2$ a $T=0$:
\[
n_e=\frac{p_F^3}{3\pi^2\hbar^3},\qquad
p_F=\hbar(3\pi^2n_e)^{1/3}.
\]
Dal bilancio dei numeri di carica:
\[
n_B=\frac AZn_e,\qquad \rho\simeq m_u n_B.
\]
Ne segue
\[
\rho\simeq\frac AZ\frac{m_u}{3\pi^2\hbar^3}p_F^3.
\]
Ogni soglia espressa in termini di $p_F$ può quindi essere convertita in una densità di massa, una volta fissato $A/Z$.
\end{frame}

% 40
\begin{frame}[t]{Atmosfera: stima del momento al quale si ionizza}
\small
Il testo utilizza una stima idrogenoide dell'energia di legame dell'elettrone più fortemente legato:
\[
E_b(Z)\simeq13.6\,Z^2\ {\rm eV}.
\]
Per elettroni non relativistici la loro energia di Fermi \emph{cinetica} è
\[
E_F^{\rm kin}\simeq\frac{p_F^2}{2m_e}.
\]
Condizione indicativa di ionizzazione/comparsa di elettroni liberi:
\[
\frac{p_F^2}{2m_e}\simeq E_b(Z)
\quad\Longrightarrow\quad p_F\simeq\sqrt{2m_eE_b(Z)}.
\]
Il modello è volutamente elementare: pressione, temperatura, shell elettroniche e fisica atomica modificano il limite reale fra atmosfera e interno ionizzato.
\end{frame}

% 41
\begin{frame}[t]{Densità di ionizzazione: legge $AZ^2$}
\small
Inseriamo $p_F^3=(2m_eE_b)^{3/2}$ nella relazione con $\rho$:
\[
\rho_{\rm ion}\simeq
\frac AZ\frac{m_u}{3\pi^2\hbar^3}
(2m_eE_b(Z))^{3/2}.
\]
Dato $E_b\propto Z^2$, segue
\[
\rho_{\rm ion}\propto \frac AZ Z^3=AZ^2.
\]
Il testo riporta la stima numerica
\[
\rho_{\rm ion}\simeq0.378\,AZ^2\ {\rm g\,cm^{-3}}.
\]
Esempi del libro: carbonio $\sim163\,{\rm g\,cm^{-3}}$; ossigeno $\sim387\,{\rm g\,cm^{-3}}$; ferro $\sim1.4\times10^4\,{\rm g\,cm^{-3}}$.
\end{frame}

% 42
\begin{frame}[t]{Quando gli elettroni diventano relativistici?}
\small
La soglia cinematica è $p_F\sim m_ec$. In unità convenzionali,
\[
\rho_{\rm rel}\simeq\frac AZ
\frac{m_um_e^3c^3}{3\pi^2\hbar^3}.
\]
La stima nel libro è
\[
\rho_{\rm rel}\simeq9.74\times10^5\frac AZ
\ {\rm g\,cm^{-3}}.
\]
Per $A/Z=2$:
\[
\rho_{\rm rel}\simeq1.95\times10^6\ {\rm g\,cm^{-3}}.
\]
\emph{Distinguere}: la relativisticità dipende da $p_F/(m_ec)$, non da $k_BT/(m_ec^2)$. Un plasma molto freddo può contenere elettroni fortemente relativistici per degenerazione.
\end{frame}

% 43
\begin{frame}[t]{Soglia di neutronizzazione: stima didattica}
\small
Il libro considera il bilancio per particelle libere
\[
p+e^-\rightarrow n+\nu_e
\]
con soglia $\mu_e\simeq m_nc^2-m_pc^2\simeq1.293\,{\rm MeV}$, trascurando intrappolamento dei neutrini.
\medskip
In regime relativistico $\mu_e\simeq\sqrt{p_F^2c^2+m_e^2c^4}$ e, come stima grossolana, $p_Fc\sim m_n c^2-m_p c^2$.
\[
\rho_{\rm neut}\sim\frac AZ\frac{m_u}{3\pi^2\hbar^3}
\left(\frac{m_nc^2-m_pc^2}{c}\right)^3.
\]
Il testo ottiene $\rho_{\rm neut}\sim3.1\times10^7\,{\rm g\,cm^{-3}}$ per $A/Z=2$.
\end{frame}

% 44
\begin{frame}[t]{Catture elettroniche reali: perché cambia la soglia?}
\small
La nana bianca contiene \emph{nuclei}, non protoni liberi. Il processo reale è, per esempio,
\[
(A,Z)+e^-\longrightarrow(A,Z-1)+\nu_e.
\]
L'energia richiesta dipende dalla differenza di masse nucleari (valore $Q$), dallo stato elettronico e dalle correzioni del reticolo.
\medskip
Quindi la stima precedente con $m_n-m_p$ serve a comprendere la scala della neutronizzazione, ma \emph{non} fornisce una soglia universale di collasso per C/O o altri nuclei.
\medskip
Si devono inoltre distinguere le catture elettroniche dalla soglia di \emph{neutron drip} tipica delle croste di stelle di neutroni.
\end{frame}

% 45
\begin{frame}[t]{Densità media e densità centrale}
\small
Per una stella con $M\sim M_\odot$ e $R\sim5000\,{\rm km}$,
\[
\bar\rho=\frac{3M}{4\pi R^3}
\sim3.8\times10^6\ {\rm g\,cm^{-3}}.
\]
La densità centrale può essere molto maggiore della media, perché $\rho(r)$ decresce verso la superficie.
\medskip
In una politropica, il rapporto $\rho_c/\bar\rho$ è fissato dall'indice $n$, non è necessariamente uno:
\[
\bar\rho=\frac{3M}{4\pi R^3}
=3\rho_c\frac{\omega_n}{\xi_1^3}.
\]
Quindi
\[
\frac{\rho_c}{\bar\rho}=\frac{\xi_1^3}{3\omega_n}.
\]
Per $n=3$ il rapporto è particolarmente elevato.
\end{frame}

% 46
\begin{frame}[t]{Spessore di uno strato sottile: stima idrostatica}
\small
Per uno strato superficiale di spessore $\Delta R\ll R$, poniamo $m(r)\simeq M$ e $r\simeq R$:
\[
\frac{dP}{dr}\simeq-\frac{GM\rho}{R^2}.
\]
Stimando $\Delta P\sim-P_b$ e $\rho\sim\rho_b$ al bordo inferiore:
\[
\frac{P_b}{\Delta R}\sim\frac{GM\rho_b}{R^2}.
\]
Ne segue
\[
\frac{\Delta R}{R}\sim\frac{P_b}{\rho_b c^2}
\left(\frac{GM}{Rc^2}\right)^{-1}.
\]
La frazione dipende dal rapporto pressione/energia e dall'inverso della compattezza $\mathcal C=GM/(Rc^2)$.
\end{frame}

% 47
\begin{frame}[t]{Spessore e pressione degli elettroni non relativistici}
\small
Per $g_e=2$:
\[
P_e^{\rm NR}=\frac{p_F^5}{15\pi^2\hbar^3m_e},\quad
n_e=\frac{p_F^3}{3\pi^2\hbar^3},\quad
\rho=\frac AZ m_un_e.
\]
Il rapporto si semplifica:
\[
\frac{P_e^{\rm NR}}{\rho c^2}
=\frac Z A\frac{p_F^2}{5m_em_uc^2}
=\frac25\frac Z A\frac{E_F^{\rm kin}}{m_uc^2}.
\]
Quindi
\[
\frac{\Delta R}{R}\sim
\frac{2}{5\mathcal C}\frac ZA\frac{E_F^{\rm kin}}{m_uc^2}.
\]
L'estensione degli strati esterni aumenta se la stella è poco compatta e se l'energia elettronica al piede dello strato è elevata.
\end{frame}

% 48
\begin{frame}[t]{Confronto con stelle di neutroni: scala degli strati}
\small
L'argomento del testo usa soltanto la stima
\[
\frac{\Delta R}{R}\sim\mathcal C^{-1}\frac{P_b}{\rho_b c^2}.
\]
La compattezza tipica è molto diversa:
\[
\mathcal C_{\rm WD}\sim10^{-4},\qquad
\mathcal C_{\rm NS}\sim0.2.
\]
A parità di $P_b/(\rho_bc^2)$, il secondo strato ha uno spessore relativo inferiore di un fattore $\sim2000$.
\medskip
Il libro fornisce stime indicative di atmosfera della nana bianca da centinaia di metri a centinaia di km a seconda della composizione. Per stelle di neutroni le scale corrispondenti sono molto più ridotte; il calcolo dettagliato richiede un'EOS e un modello atmosferico realistici.
\end{frame}

\section{Effetti termici e cristallizzazione}
\subsection{§5.7: gas degenere e plasma coulombiano}

% 49
\begin{frame}[t]{Due scale energetiche distinte: $k_BT$ e $E_F$}
\small
L'approssimazione di gas completamente degenere è valida se
\[
k_BT\ll E_F^{\rm kin}
\]
(con le opportune generalizzazioni per elettroni relativistici).
\medskip
Conversione utile:
\[
k_B\simeq8.617\times10^{-5}\ {\rm eV\,K^{-1}},\quad
1\,{\rm MeV}\leftrightarrow1.16\times10^{10}\,{\rm K}.
\]
Per esempio $T=10^4\,{\rm K}$ corrisponde a $k_BT\simeq0.86\,{\rm eV}$.
\medskip
La condizione $k_BT\ll m_ec^2$ \emph{non} significa che gli elettroni del gas siano non relativistici: è il rapporto $p_F/(m_ec)$ a decidere la cinematica.
\end{frame}

% 50
\begin{frame}[t]{Temperatura di Fermi per elettroni non relativistici}
\small
Con
\[
p_F=\hbar(3\pi^2n_e)^{1/3},\qquad
E_F^{\rm kin}=\frac{p_F^2}{2m_e},
\]
definiamo
\[
k_BT_F\equiv E_F^{\rm kin}
=\frac{\hbar^2(3\pi^2)^{2/3}}{2m_e}n_e^{2/3}.
\]
Sostituendo $n_e=\rho/(\mu_em_u)$:
\[
T_F=\frac{\hbar^2(3\pi^2)^{2/3}}{2m_ek_B}
\left(\frac\rho{\mu_em_u}\right)^{2/3}.
\]
L'energia di Fermi cresce con la densità come $\rho^{2/3}$; anche a temperature macroscopiche elevate l'interno stellare può quindi essere fortemente degenere.
\end{frame}

% 51
\begin{frame}[t]{Densità critica dalla condizione $E_F^{\rm kin}=k_BT$}
\small
Uguagliamo le scale:
\[
\frac{\hbar^2(3\pi^2 n_e)^{2/3}}{2m_e}=k_BT.
\]
Risolvendo per $n_e$:
\[
n_{e,\rm crit}=\frac1{3\pi^2}
\left(\frac{2m_ek_BT}{\hbar^2}\right)^{3/2}.
\]
Per la densità di massa:
\[
\rho_{\rm crit}=\mu_em_u n_{e,\rm crit}
\propto\mu_eT^{3/2}.
\]
Il libro riporta
\[
\rho_{\rm crit}\simeq0.0121\,\frac{A}{2Z}
\left(\frac{T}{10^4\,{\rm K}}\right)^{3/2}
{\rm g\,cm^{-3}}.
\]
È una stima di degenerazione, non la densità di ionizzazione o di cristallizzazione.
\end{frame}

% 52
\begin{frame}[t]{Secondo criterio: uguaglianza delle pressioni}
\small
Nel limite classico $P_{\rm cl}=n_ek_BT$; per gli elettroni degeneri non relativistici:
\[
P_{\rm deg}=\frac{\hbar^2(3\pi^2)^{2/3}}{5m_e}n_e^{5/3}
=\frac25n_eE_F^{\rm kin}.
\]
La condizione $P_{\rm cl}=P_{\rm deg}$ fornisce
\[
k_BT=\frac25E_F^{\rm kin}.
\]
A $T$ fissata richiede densità maggiore rispetto a $k_BT=E_F^{\rm kin}$:
\[
\frac{\rho_{\rm crit}^{(P)}}{\rho_{\rm crit}^{(E)}}
=\left(\frac52\right)^{3/2}\simeq3.95.
\]
La differenza è un coefficiente numerico; in entrambi i casi $\rho_{\rm crit}\propto T^{3/2}$.
\end{frame}

% 53
\begin{frame}[t]{Oceano liquido e cristallo coulombiano}
\small
La degenerazione riguarda gli \emph{elettroni}; il reticolo dei \emph{nuclei} può fondere o cristallizzare in funzione del rapporto fra energia coulombiana e temperatura.
\medskip
Per ioni di carica $Ze$ separati da un raggio medio $a_i$:
\[
\Gamma_C\equiv\frac{Z^2e^2}{a_ik_BT}
\quad(e^2=\alpha\hbar c).
\]
\begin{itemize}
\item $\Gamma_C$ piccolo: agitazione termica prevale, fase liquida / plasma ionico.
\item $\Gamma_C$ grande: le correlazioni coulombiane favoriscono un solido ordinato.
\end{itemize}
Per un plasma one-component classico il libro usa $\Gamma_{C,m}\simeq175$ e, per il caso adottato di carbonio, una soglia vicina a $178$.
\end{frame}

% 54
\begin{frame}[t]{Temperatura di fusione: derivazione dalla cella ionica}
\small
Il raggio della cella è fissato da $n_i=\rho/(Am_u)$:
\[
a_i=\left(\frac{3Am_u}{4\pi\rho}\right)^{1/3}.
\]
All'equilibrio di fusione $\Gamma_C=\Gamma_{C,m}$; quindi
\[
k_BT_m=\frac{Z^2\alpha\hbar c}{\Gamma_{C,m}a_i}.
\]
Sostituendo $a_i$:
\[
T_m=\frac{Z^2\alpha\hbar c}{\Gamma_{C,m}k_B}
\left(\frac{4\pi\rho}{3Am_u}\right)^{1/3}.
\]
Per composizione fissata:
\[
T_m\propto Z^2A^{-1/3}\rho^{1/3}.
\]
L'aumento della densità favorisce la cristallizzazione a temperatura più alta.
\end{frame}

% 55
\begin{frame}[t]{Cosa significa $\Gamma_C\simeq178$ per una nana bianca?}
\small
Il parametro Coulomb misura la competizione fra ordine elettrostatico e agitazione termica; non è l'indice politropico $\Gamma$ del §5.2.
\medskip
Per $T<T_m(\rho)$ il materiale ionico tende a essere solido; per $T>T_m(\rho)$ può formare l'\emph{ocean} liquido sopra gli strati cristallizzati.
\medskip
Le variabili termiche della nana bianca non sono uniformi: la temperatura interna può essere molto diversa dalla temperatura efficace superficiale $T_{\rm eff}$.
\medskip
Il risultato illustrato nel libro è soprattutto qualitativo: la presenza di un oceano e il fronte di cristallizzazione influenzano il raffreddamento e l'evoluzione termica, pur lasciando la pressione di degenerazione elettronica come sostegno principale.
\end{frame}

\section{Osservazioni delle nane bianche}
\subsection{§5.8: massa, raggio e modelli atmosferici}

% 56
\begin{frame}[t]{Produzione stellare e tipi di nane bianche}
\small
Il libro descrive le nane bianche come resti di stelle che espellono gli strati esterni, lasciando soprattutto core di C/O in molti percorsi evolutivi.
\medskip
Una stima semplificata di durata della sequenza principale usa
\[
L\propto M^\alpha,\quad\alpha\sim4,\qquad
\tau_{\rm MS}\sim\frac ML\propto M^{1-\alpha}\sim M^{-3}.
\]
Stelle più massicce consumano prima il combustibile.
\medskip
\emph{Precisazione}: la conclusione del libro che quasi tutte le nane bianche presenti debbano essere C/O è una semplificazione. Le nane bianche di elio a bassa massa possono formarsi tramite evoluzione binaria con perdita dell'involucro, e non vanno escluse in generale.
\end{frame}

% 57
\begin{frame}[t]{Classi spettrali: che cosa osserviamo?}
\small
I tipi spettroscopici citati nella Tabella 5.2 sono identificati dalle righe osservate:
\[
\begin{array}{c|l}
\text{classe}&\text{caratteristica}\\\hline
\mathrm{DA}&\text{righe dell'idrogeno (Balmer)}\\
\mathrm{DB}&\text{righe dell'elio}\\
\mathrm{DC}&\text{continuo, senza righe evidenti}\\
\mathrm{DQ}&\text{tracce spettrali del carbonio}\\
\mathrm{DZ}&\text{righe di elementi pesanti}
\end{array}
\]
La composizione atmosferica altera i profili di riga, la temperatura efficace dedotta e la gravità superficiale stimata. Una classificazione atmosferica non identifica direttamente la composizione del \emph{core}.
\end{frame}

% 58
\begin{frame}[t]{Energia della radiazione di corpo nero}
\small
Per fotoni in equilibrio termico, con due polarizzazioni:
\[
f(p)=\frac1{e^{pc/(k_BT)}-1},\quad g_\gamma=2.
\]
La densità di energia è
\[
u_\gamma=g_\gamma\int\frac{d^3p}{(2\pi\hbar)^3}\,pc\,f(p).
\]
Nel sistema $\hbar=c=k_B=1$:
\[
u_\gamma=\frac1{\pi^2}\int_0^\infty\frac{p^3dp}{e^{p/T}-1}
=\frac{\pi^2}{15}T^4,
\]
perché $\int_0^\infty x^3/(e^x-1)\,dx=\pi^4/15$.
\end{frame}

% 59
\begin{frame}[t]{Densità di energia e flusso: perché compare $1/4$?}
\small
Il flusso radiativo uscente per distribuzione isotropa nel sistema comovente è l'integrale sull'emisfero:
\[
F=\frac{c\,u_\gamma}{4\pi}\int_0^{2\pi}d\phi
\int_0^{\pi/2}\sin\theta\cos\theta\,d\theta.
\]
Calcolando:
\[
\int_0^{\pi/2}\sin\theta\cos\theta\,d\theta=\frac12,
\qquad \int_0^{2\pi}d\phi=2\pi.
\]
Perciò
\[
F=\frac{cu_\gamma}{4}=\sigma_{\rm SB}T^4,
\qquad \sigma_{\rm SB}=\frac{\pi^2 k_B^4}{60\hbar^3c^2}.
\]
Il fattore $1/4$ deriva dal peso angolare della direzione di emissione, non da una riduzione arbitraria.
\end{frame}

% 60
\begin{frame}[t]{Legge di Stefan--Boltzmann e misura del raggio}
\small
Per una nana bianca approssimata come corpo nero di raggio $R$:
\[
L=4\pi R^2\sigma_{\rm SB}T_{\rm eff}^4.
\]
L'osservatore sulla Terra misura il flusso bolometrico a distanza $D$:
\[
F_\oplus=\frac{L}{4\pi D^2}
=\sigma_{\rm SB}T_{\rm eff}^4\frac{R^2}{D^2}.
\]
Invertendo:
\[
R=D\sqrt{\frac{F_\oplus}{\sigma_{\rm SB}T_{\rm eff}^4}}.
\]
Serve una distanza misurata (per esempio da parallasse); l'analisi atmosferica reale usa il flusso integrato nello spettro anziché un corpo nero perfetto.
\end{frame}

% 61
\begin{frame}[t]{Parallasse, gravità superficiale e redshift}
\small
Tre osservabili complementari:
\begin{itemize}
\item \textbf{Parallasse}: determina $D$, quindi $R$ dal flusso e da $T_{\rm eff}$.
\item \textbf{Profili spettrali}: stimano la gravità $g\simeq GM/R^2$, spesso espressa come $\log_{10}g$ in unità cgs.
\item \textbf{Redshift gravitazionale}: nel limite debole
\[
z\simeq\frac{GM}{Rc^2},\qquad
1+z=\left(1-\frac{2GM}{Rc^2}\right)^{-1/2}
\]
per fotoni emessi da una superficie statica sferica e ricevuti all'infinito.
\end{itemize}
Se $M$ è noto dalla dinamica binaria, $g$ o $z$ permettono di determinare $R$ con ipotesi differenti.
\end{frame}

% 62
\begin{frame}[t]{Perché la massa derivata da un modello non è un test indipendente?}
\small
Supponiamo che un catalogo riporti una massa inferita da $T_{\rm eff}$ e $\log g$ \emph{adottando già} una relazione teorica $M(R)$.
\medskip
Confrontare la massa così dedotta con la stessa relazione $M(R)$ sarebbe circolare.
\medskip
Per una verifica non tautologica occorrono vincoli largamente indipendenti, ad esempio:
\begin{enumerate}
\item massa dinamica da orbite binarie o lente gravitazionale;
\item raggio da parallasse, flusso e temperatura con modello atmosferico;
\item oppure redshift gravitazionale insieme a una massa dinamica.
\end{enumerate}
Il libro privilegia alcuni sistemi vicini con misure indipendenti di massa e raggio.
\end{frame}

% 63
\begin{frame}[t]{Tabella 5.3: quattro nane bianche vicine}
\small
\begin{center}
\begin{tabular}{lrrr}
\hline
Oggetto & $T_{\rm eff}$ (K) & $R/R_\odot$ & $M/M_\odot$\\\hline
Procyon B & $7740\pm50$ & $0.01232\pm0.00032$ & $0.592\pm0.006$\\
Sirius B & $25369\pm46$ & $0.008098\pm0.000046$ & $1.018\pm0.011$\\
Stein 2051 B & $7122\pm181$ & $0.0114\pm0.0004$ & $0.675\pm0.051$\\
40 Eridani B & $17200\pm110$ & $0.01308\pm0.00020$ & $0.573\pm0.018$\\\hline
\end{tabular}
\end{center}
\medskip
Valori riportati nel libro (Bond et al. 2017b): sono dati storici del confronto proposto nel capitolo, non un catalogo osservativo aggiornato.
\medskip
Notare l'andamento generale: a masse maggiori corrispondono raggi inferiori.
\end{frame}

% 64
\begin{frame}[t]{Figura 5.1: diagramma massa--raggio delle nane bianche}
\centering
\begin{figure}[!ht]\includegraphics[trim={0cm 0cm 0 0},clip,keepaspectratio,width=0.95\textwidth]{Schaffner/MR_WD}\label{fig:Schaffner_MR_WD}\end{figure}
\end{frame}

% 65
\begin{frame}[t]{Interpretazione della Figura 5.1}
\small
La figura confronta quattro sistemi vicini con famiglie teoriche massa--raggio:
\begin{itemize}
\item Core di C/O ($Z/A\simeq1/2$) e core di Fe ($Z/A\simeq0.464$) producono curve differenti.
\item Le correzioni termiche e i modelli di atmosfera spostano il raggio dedotto a massa fissata.
\item Il parametro $q_H=M_H/M_{\rm WD}$ indica la frazione massica di uno strato superficiale di H (nel grafico: atmosfere spesse o sottili).
\end{itemize}
I dati citati dal libro sono in accordo, entro il modello adottato, con core C/O; non richiedono una nana bianca di ferro per spiegarne massa e raggio.
\medskip
Il confronto usa modelli atmosferici, non soltanto le politropiche ideali $n=3/2$ o $3$.
\end{frame}

% 66
\begin{frame}[t]{Effetti termici sui raggi: attenzione all'interpretazione}
\small
Nella figura, l'elevata $T_{\rm eff}$ di Sirius B e 40 Eridani B è correlata a modelli con raggi maggiori rispetto a sequenze più fredde.
\medskip
Non bisogna però dedurre che tutta la stella abbia temperatura uniforme $T_{\rm eff}$: il nucleo è più caldo dell'atmosfera.
\medskip
Il raggio dipende da composizione del core, spessore degli strati superficiali e struttura termica. La pressione elettronica rimane largamente degenere nell'interno; le correzioni termiche possono comunque essere importanti per le regioni esterne.
\end{frame}

\section{Sintesi e collegamenti}

% 67
\begin{frame}[t]{Ch5: catena di derivazioni da ricordare}
\small
\[
P=K'\rho^\Gamma,\qquad \Gamma=1+\frac1n
\]
\[
\frac1{r^2}\frac d{dr}\left(\frac{r^2}{\rho}\frac{dP}{dr}\right)
=-4\pi G\rho
\]
\[
\rho=\rho_c\theta^n,\quad r=a\xi
\quad\Longrightarrow\quad
\frac1{\xi^2}(\xi^2\theta')'=-\theta^n
\]
\[
M=4\pi a^3\rho_c\omega_n,\qquad R=a\xi_1
\]
\[
n=3\ (\Gamma=4/3)\quad\Longrightarrow\quad
M_{\rm Ch}\propto\frac{m_{\rm Pl}^3}{(\mu_e m_u)^2}.
\]
La massa di Chandrasekhar collega statistica quantistica, relatività speciale e gravità newtoniana.
\end{frame}

% 68
\begin{frame}[t]{Dal capitolo 5 ai capitoli sulle stelle di neutroni}
\small
\begin{itemize}
\item La pressione di Fermi degli elettroni produce la sequenza delle nane bianche; per stelle di neutroni entra invece la microfisica nucleare e la relatività generale.
\item La politropica $\Gamma=4/3$ genera una massa limite nel modello newtoniano, analoga per scala all'argomento di Landau del §4.2.
\item Lane--Emden mostra il ruolo di struttura e condizioni al bordo; le TOV ne estendono il problema in relatività generale.
\item Le correzioni coulombiane, la composizione e la temperatura sono esempi di ciò che una EOS puramente ideale non descrive.
\end{itemize}
\medskip
\emph{Riferimento}: capitolo 5 di \emph{Compact Star Physics}, sezioni 5.1--5.8. I chiarimenti su segni, soglie di cattura elettronica ed evoluzione binaria sono esplicitati come precisazioni degli appunti.
\end{frame}

